\section{Critical moments and susceptibility}
\label{mom:section}

We first turn the local maps into a lower bound on the rate of growth of
the susceptibility.  Under the temporary assumption \eqref{eq:s09},
\cref{amp:efficient-maps} supplies maps on a fixed orbit $O$ with
\[
 U=\log(1/u),\qquad m=\log M,\qquad H=\log(1/\norm K),
 \qquad \frac{U+m}{H}\longrightarrow0.
\]
The calculations with susceptibilities and second moments take place
at $p<\pc$, where the preliminary estimates ensure their finiteness.

\begin{lemma}[Maximal susceptibility in an intrinsic ball]
\label{mom:maximal-potential}
For a deterministic induced domain $D$, put
$f_D(v)=\E_p\abs{C_v^D}$ for $v\in D$, and $f_D(v)=0$ otherwise.
Let $\mathcal B^D(v,M)$ be the vertices reachable from $v$ by an open
path in $D$ of length at most the nonnegative integer $M$.  If $v\in D$
and $0<p<\pc$, then
\begin{equation}
 \E_p\max_{z\in\mathcal B^D(v,M)}f_D(z)
 \le \left(1+\frac{1-p}{p}M\right)f_D(v).
 \label{eq:s16}
\end{equation}
\end{lemma}

\begin{proof}
First let $D$ be finite.  Select a maximizer by a fixed tie rule.  For
each $w\in D$, take a percolation configuration $\omega$ and expose the
exact vertex set $B=C_w^D(\omega)$.  Replace every edge with at least one
endpoint in $B$ by an independent Bernoulli($p$) bit, leaving all other
edges unchanged; call the resulting configuration $\widetilde\omega$.
Conditional on $B$, the unchanged exterior edges have product law,
and all replaced bits are fresh.  Consequently $\widetilde\omega$ has
product law independent of $B$.  If $Z$ is the selected maximizer in
$\mathcal B^D(v,M;\widetilde\omega)$, this independence gives
\[
 \sum_{w\in D}\Prob(Z\in B)
 =\sum_{z\in D}\Prob(Z=z)\sum_{w\in D}T_p^D(w,z)
 =\E_p\max_{z\in\mathcal B^D(v,M)}f_D(z).
\]
Here a separate coupling is used for each $w$; only the marginal law of
$Z$ needs to be the same in these couplings.

If $v\notin B$ and $Z\in B$, take the first edge entering $B$ along
a simple $\widetilde\omega$-open path from $v$ to $Z$ of length at most
$M$.  This edge was closed in $\omega$, and its exterior endpoint was
reachable from $v$ in $D\setminus B$ by at most $M-1$ open edges of
$\omega$.  Let $N_w(\omega)$ count boundary edges with this property,
setting it to zero when $v\in B$.  Then
\[
 \Prob(Z\in B)\le T_p^D(v,w)+\E_p N_w.
\]
To bound the last expectation, open one counted, distinguished edge.
In the new configuration it is pivotal for $v\leftrightarrow w$:
it is the only open edge entering the old cluster $B$.  Moreover it is
among the first $M$ bridges separating $v$ from $w$, in their order
from $v$, because a simple path reaches it using at most $M$ edges.
For a fixed output and a distinguished edge, closing that edge
recovers the input.  There are at most $M$ possible distinguished
edges in each connected output.  Comparing the two possible values
of that one bit therefore gives
\[
 \E_p N_w\le \frac{1-p}{p}M\,T_p^D(v,w).
\]
Summation over $w$ proves \eqref{eq:s16}.  For an infinite domain,
exhaust $D$ by finite induced domains containing $v$.  Their
susceptibilities and intrinsic balls increase to the corresponding
objects in $D$.  The maximum also converges monotonically: the
intrinsic $M$-ball is contained in a deterministic finite graph ball,
and every witnessing path is finite.  Monotone convergence proves
the assertion.  The case $M=0$ is immediate.
\end{proof}

\begin{lemma}[A local map in a deleted domain]
\label{mom:map-potential}
Fix one of the maps above, with critical success probability $u>0$,
cap $M\ge2$, and $H>1$.  Run its fixed rule at parameter $p$.  On a
map-dependent interval $p_0\le p<\pc$, with $p_0\ge\pc/2$, its success
probability $u_p$ and conditional endpoint kernel $K_p$ satisfy
\[
 u_p\ge u/2,\qquad \norm{K_p}\le e^{-(H-1)}.
\]
For a deterministic set $A\subseteq V$, put
$F_A=f_{V\setminus A}/\chi$, with value zero on $A$.  There are
constants $Z,C_{\mathrm{map}}$, uniform on this interval, such that
$2\le Z\le CM/u$ and, for every $v\in O$,
\begin{equation}
 (K_pF_A)(v)\le ZF_A(v)+C_{\mathrm{map}}h_A(v),
 \qquad
 h_A(v)=\sum_{\substack{a\in A\\\dist(v,a)\le M}}
                    \sum_{b\sim a}F_A(b).
 \label{eq:s17}
\end{equation}
In the action of $K_p$, functions are restricted to $O$.  The constant
$C$ depends only on the graph.
\end{lemma}

\begin{proof}
The unconditional endpoint kernel has finite range and continuous
entries as a function of $p$: its rule queries only finitely many
edge bits, and its auxiliary randomization has a fixed law.
Covariance makes its row sums constant on $O$.  The absolute
difference of two such kernels is invariant on $O\times O$, so
\eqref{eq:s02} gives equal row and column sums.  The finite range and
the row--column norm bound therefore give operator-norm continuity.
Dividing by the positive scalar $u_p\to u$ proves the assertions
after shortening the interval.  The kernel $K_p$ is stochastic and
invariant, and hence has column sums one, although it need not be
symmetric.  The cap and connection conditions are properties of the
accepted local patterns, so every success at $p$ still has a simple
open path of length at most $M$ to its endpoint.

Choose such a path on each success.  If it avoids $A$, its endpoint
lies in the intrinsic $M$-ball in $V\setminus A$.  Thus
\cref{mom:maximal-potential} bounds its unconditional contribution to
the endpoint potential by
$\bigl(1+(1-p)M/p\bigr)F_A(v)$; this case is empty if $v\in A$.
Endpoints in $A$ contribute zero.  In the remaining case let
$a\sim b$ be the last exit of the chosen path from $A$.  Its terminal
segment from $b$ to the endpoint $y$ lies in $V\setminus A$ and has
length $l\le M$.  Positive association with the event that this
deterministic segment is open yields
\[
 f_{V\setminus A}(b)
 \ge \E_p\bigl[\abs{C_y^{V\setminus A}};
                       \text{the segment is open}\bigr]
 \ge p^l f_{V\setminus A}(y).
\]
Therefore $F_A(y)\le p^{-M}F_A(b)$.  Since
$\dist(v,a)\le M$, summing over all possible last exits bounds this
contribution by $p^{-M}h_A(v)$.  Divide the two bounds by $u_p$.
Because $p\ge\pc/2$ and $u_p\ge u/2$, the first coefficient can be
enlarged to a constant $Z\in[2,CM/u]$, and the second to a finite
constant depending on the fixed map.  This proves \eqref{eq:s17}.
\end{proof}

\begin{proposition}[Critical fractional moments]
\label{mom:critical-moments}
For every $0<\alpha<1/2$,
\begin{equation}
 \sup_{v\in V}\E_{\pc}s_v^\alpha<\infty.
 \label{eq:s21}
\end{equation}
\end{proposition}

\begin{proof}
Assume \eqref{eq:s09}, and define
\[
 D(p)=\chi^{-2}\left[
       \sum_{b\sim o}\E_p(s_os_b;\ C_o\ne C_b)\right].
\]
By \eqref{eq:s04}, $(1-p)\chi'=\chi^2D(p)$.  Our first goal is to
prove, for every $\gamma>0$, that $D(p)\ge c_\gamma\chi^{-\gamma}$
when $p$ is sufficiently close to $\pc$.  We will compare two bounds
on the exterior potential reached by repeated endpoint-map steps:
\eqref{eq:s17} bounds it above in terms of $D(p)$, while the small
norm of $K_p$ bounds the contribution of endpoints in the root's
own cluster and hence leaves a positive exterior contribution.

Fix a map and its interval as in \cref{mom:map-potential}.  For
$x\in O$, expose $A=C_x$ in an independent $p$-field, and use the
possibly unnormalized expectation
\[
 \widehat\E[G(A)]=\E_p\left[\frac{s_x}{\chi}G(C_x)\right].
\]
For each fixed vertex $a$, exposure of the exact cluster $C_a$
leaves fresh exterior edges and hence gives
\begin{align*}
 \widehat\E\left[\one_{\{a\in A\}}
                         \sum_{b\sim a}F_A(b)\right]
 &=\chi^{-2}\sum_{b\sim a}
      \E_p\left[s_a\one_{\{x\in C_a\}}f_{V\setminus C_a}(b)\right]
 \\
 &\le \chi^{-2}\sum_{b\sim a}
                 \E_p(s_as_b;\ C_a\ne C_b)
 \le c_\sigma D(p),
\end{align*}
where $c_\sigma=(\min_i\sigma_i)^{-1}$.  The last inequality follows
because the preceding nonnegative expression is constant on each
vertex type and its type average is $D(p)$.  Bounded degree now gives
$\widehat\E h_A(v)\le C_{\mathrm{map}}D(p)$ uniformly in $v$.
Iterate \eqref{eq:s17}, use stochasticity of $K_p$, and note that
$F_A(x)=0$.  Since $Z\ge2$, the geometric sum gives
\begin{equation}
 \widehat\E(K_p^jF_A)(x)
       \le C_{\mathrm{map}} Z^jD(p),\qquad j\ge1.
 \label{eq:s18}
\end{equation}

We next give a lower bound for the same quantity.  Fresh exterior
exposure identifies it as
\[
 \widehat\E(K_p^jF_A)(x)
 =\chi^{-2}\sum_{v\in O}K_p^j(x,v)
                  \E_p(s_xs_v;\ C_x\ne C_v).
\]
Without the distinctness condition, positive association and the
uniform near-critical comparison of susceptibilities give a lower
bound $c_0>0$ for this normalized sum.  To control what is removed,
write $I=C_x\cap O$, which is nonempty.  Mass transport restricted to
$O$ gives the exact identity
\begin{equation}
 \sum_{v\in O}K_p^j(x,v)\E_p(s_x^2;\ v\in C_x)
 =\E_p\left[\frac{s_x^2}{\abs I}
                  \ip{\one_I}{K_p^j\one_I}\right]
 \le \norm{K_p}^{j}\E_p s_x^2
 \le C\norm{K_p}^{j}\chi^3.
 \label{mom:diagonal-loss}
\end{equation}
For completeness, the transport sends, from $z\in O$ to each
$y\in C_z\cap O$, the mass
\[
 \frac{s_z^2}{\abs{C_z\cap O}}
       \sum_{v\in C_z\cap O}K_p^j(z,v).
\]
Its outgoing and incoming expectations are the two sides of the
identity.  The norm bound uses
$\norm{\one_I}_2^2=\abs I$; it does not require symmetry of $K_p$.
Finally, the three-arm estimate \eqref{eq:s01} gives
$\E_p s_x^2\le\chi_{\max}^3\le C\chi^3$.

Choose $C_*>1$ large enough that
\[
 j=\left\lceil\frac{\log(C_*\chi)}{H-1}\right\rceil
\]
makes the normalized loss in \eqref{mom:diagonal-loss} at most
$c_0/2$.  Combining this lower bound with \eqref{eq:s18} and
$\log Z\le U+m+C'$ gives
\begin{equation}
 D(p)\ge c_{\mathrm{map}}
               \chi^{-(U+m+C')/(H-1)}.
 \label{eq:s19}
\end{equation}
Here $C'$ is graph-dependent; the prefactor may depend on the fixed
map.  Since $U+m\ge\log2$, the maps with $(U+m)/H\to0$ have
$H\to\infty$.  Thus, for any $0<\gamma<1$, we can fix a map for
which the exponent in \eqref{eq:s19} is at most $\gamma$.
The susceptibility identity \eqref{eq:s04} then gives, almost
everywhere on its interval,
$\chi'\ge c\chi^{2-\gamma}$.  In particular
\[
 \frac{d}{dp}\chi^{-(1-\gamma)}
       \le -(1-\gamma)c.
\]
Integrating from $p$ to $r<\pc$, discarding the nonnegative value at
$r$, and then letting $r\uparrow\pc$ yields
\begin{equation}
 \chi(p)\le C_\gamma(\pc-p)^{-1/(1-\gamma)}.
 \label{eq:s20}
\end{equation}
Local absolute continuity, established with \eqref{eq:s04},
justifies this differential calculation.

We transfer this bound to a critical tail by an explicit comparison
of exploration transcripts.  For a large integer $N$, put
$p=\pc-N^{-1/2}$.  From any vertex $v$, explore until the cluster
is exhausted or $N$ vertices have been discovered.  Use a fixed
nonretesting exploration order.  There are at most $dN$ fresh edge
queries, and the transcript decides the event $s_v\ge N$.
If $L_N$ is the likelihood ratio of its law at $\pc$ relative to
its law at $p$, each active query contributes either $\pc/p$ or
$(1-\pc)/(1-p)$.  For fixed $b>1$, its conditional $b$th-moment
factor under $p$ is
\[
 p(\pc/p)^b+(1-p)((1-\pc)/(1-p))^b
       \le 1+C_b(\pc-p)^2.
\]
Indeed this expression is twice continuously differentiable in a
fixed neighborhood of the interior point $\pc$, equals one there,
and has first derivative zero there.  Pad the stopped exploration
with inactive factors equal to one.  Successive conditioning for
the deterministic $dN$ steps gives
\[
 \E_p L_N^b\le(1+C_b/N)^{dN}\le e^{dC_b}.
\]
H\"older's inequality and Markov's inequality therefore imply,
uniformly in $v$,
\begin{equation}
 \Prob_{\pc}(s_v\ge N)
 =\E_p[L_N\one_{\{s_v\ge N\}}]
 \le C_b'\left(\frac{\chi_v(p)}{N}\right)^{1-1/b}.
 \label{mom:transcript-comparison}
\end{equation}
Choose $\gamma$ small and then $b$ large enough that
\[
 \theta=(1-1/b)\left(1-\frac1{2(1-\gamma)}\right)>\alpha,
\]
where $\alpha$ is the exponent in \eqref{eq:s09}.  For all
sufficiently large $N$, \eqref{eq:s20} applies to the chosen $p$;
the uniform type comparison and
\eqref{mom:transcript-comparison} give
$\Prob_{\pc}(s_v\ge N)\le C N^{-\theta}$.
Summing
$((N+1)^\alpha-N^\alpha)\Prob_{\pc}(s_v>N)$ proves a uniform finite
$\alpha$th moment, contrary to \eqref{eq:s09}.  This establishes
\eqref{eq:s21} for every $0<\alpha<1/2$.
\end{proof}

The next argument uses these moments without the temporary
divergence assumption. The conversion from a critical volume tail to
a susceptibility estimate is the principle established in
\cite[Theorem~1.1]{Hutchcroft2020Exponents} for transitive graphs.
We prove the version needed here with constants uniform over the
finitely many vertex types. We include the covariance form of the
O'Donnell--Saks--Schramm--Servedio inequality
\cite[Section~3.3, Equation~(6)]{ODonnellSaksSchrammServedio2005},
in the form needed for all vertex types and all finite domains.

\begin{lemma}[A decision-tree covariance bound]
\label{mom:decision-tree}
Let $X$ have a finite product law, let $0\le g\le1$, and let an
algorithm determine $g(X)$ by querying coordinates without
repetition.  Write $\delta_i$ for the probability that it queries
coordinate $i$, and let $X^{(i)}$ replace just that coordinate by an
independent copy.  For every real function $f$ on the product space,
\begin{equation}
 \abs{\Cov(f(X),g(X))}
 \le\sum_i\delta_i\E\abs{f(X)-f(X^{(i)})}.
 \label{mom:covariance}
\end{equation}
Independent randomization of the algorithm is allowed.
\end{lemma}

\begin{proof}
Take an independent product copy $Y$.  Run the algorithm on $X$,
and let $Z^t$ replace the first $t$ queried coordinates of $X$ by
their values in $Y$.  Include the algorithm's independent seed in
each transcript.  Conditional on any pre-query transcript, the
unqueried coordinates of $X$ retain their independent original
marginals, since that transcript constrains only queried bits.
The replaced coordinates of $Y$ are independent of this transcript
and of those unqueried bits.  Thus the pre-query hybrid has the
original product law conditional on the transcript.  If the next
index is $i$, that index is transcript-measurable and $Y_i$ is still
fresh.  Consequently
\[
 \E\bigl[\abs{f(Z^{t-1})-f(Z^t)}
                     \mid\text{pre-query transcript}\bigr]
       =\E\abs{f(X)-f(X^{(i)})}.
\]
At the terminal transcript, the same reasoning gives product law
for $Z^T$ conditional on the entire transcript.  Since the
transcript determines $g(X)$, it follows that
\[
 \Cov(f(X),g(X))=\E[(f(X)-f(Z^T))g(X)].
\]
Telescope and use $0\le g\le1$ before conditioning on partial
transcripts.  The preceding conditional-increment identity then
proves \eqref{mom:covariance}, because each coordinate is queried
at most once.  Intermediate hybrids need not be independent of
$g(X)$.
\end{proof}

\begin{proposition}[An almost linear susceptibility bound]
\label{mom:susceptibility}
For every $\eta>0$ there is $C_\eta<\infty$ such that
\begin{equation}
 \chi_{\max}(p)\le C_\eta(\pc-p)^{-1-\eta},
 \qquad 0\le p<\pc.
 \label{eq:s22}
\end{equation}
\end{proposition}

\begin{proof}
Fix $0<a<1/2$.  In a finite induced domain, fix a vertex $v$ and
an integer $n\ge2$.  Let $f$ indicate $\abs{C_v}\ge n$.
Independently mark every vertex with probability $1/n$, and let $g$
indicate that $C_v$ contains a mark.  An algorithm deciding $g$
first reveals all marks and then explores the clusters of marked
vertices.  It queries an edge only if at least one endpoint belongs
to a cluster containing a mark.  By \cref{mom:critical-moments},
restriction coupling, and monotonicity in $p\le\pc$, each vertex
has this property with probability at most
\[
 \E_p[1-(1-1/n)^{s_v}]
 \le\E_p\min(1,s_v/n)
 \le n^{-a}\E_p s_v^a\le C_a n^{-a}.
\]
Thus every edge has revealment at most $2C_a n^{-a}$, uniformly in
the domain, the root, and $p\le\pc$.  Mark coordinates have zero
influence on $f$.  For an increasing Boolean edge function, its
absolute resampling influence is exactly
$2p(1-p)$ times the probability that the edge is pivotal.
Writing $F(p)=\Prob_p(s_v\ge n)$, finite-product differentiation
and \cref{mom:decision-tree} give
\[
 \Cov(f,g)\le C_a n^{-a}\,2p(1-p)F'(p).
\]
On the other hand, conditioning on the percolation configuration
gives
\[
 \Cov(f,g)
 \ge\bigl(1-(1-1/n)^n-C_a n^{-a}\bigr)F(p).
\]
For all sufficiently large $n$, the factor in parentheses is
bounded below by a positive constant.  Since $2p(1-p)\le1/2$,
we obtain $F'(p)\ge c_a n^a F(p)$ for $0<p\le\pc$.
Multiplying by $\exp(-c_an^ap)$ shows that
$\exp(-c_an^ap)F(p)$ is nondecreasing.  Integrating up to $\pc$
therefore gives, including the case $F\equiv0$,
\begin{equation}
 \Prob_p(s_v\ge n)
 \le C_a n^{-a}\exp\{-c_a(\pc-p)n^a\}.
 \label{mom:stretched-tail}
\end{equation}
We used the critical moment bound to estimate $F(\pc)$.
Continuity gives the assertion also at $p=0$.  Increasing finite
domain exhaustion is legitimate because the event $s_v\ge n$ has
a finite open witness; all constants are independent of the
domain.  Enlarge the constant to cover the finitely many smaller
values of $n$.

Put $\delta=\pc-p$.  The function
$x^{-a}e^{-c_a\delta x^a}$ is decreasing on $(0,\infty)$.
The substitution $t=c_a\delta x^a$ in its integral gives
\[
 \sum_{n\ge1}n^{-a}e^{-c_a\delta n^a}
 \le C_a\delta^{-(1-a)/a}.
\]
Summing \eqref{mom:stretched-tail} therefore proves the same
bound for $\chi_{\max}$.  Finally take $a=1/(2+\eta)$, for which
$(1-a)/a=1+\eta$, to obtain \eqref{eq:s22}.
\end{proof}
